%!TEX root = paper.tex
\subsection{Dataset Construction and Artifact Extraction}\label{sec:method:dataset}

We construct datasets from multiple sources to reduce bias.  The first dataset contains popular PyPI packages selected by download counts or dependency centrality.  The second contains randomly sampled PyPI packages across popularity and age.  The third contains packages reported as malicious or suspicious by public advisories, security vendors, or package-index takedowns.  The fourth contains Python application bundles and release artifacts such as wheels, archives, container layers, and PyInstaller-style outputs.  The fifth contains runtime-testing inputs derived from CPython tests, real bytecode artifacts, and generated bytecode variants.

For each package or artifact, {\tech} records the package name and version, distribution type, file inventory, source and bytecode paths, bytecode magic number, inferred Python version, hashes, and import or load behavior when such behavior can be determined statically.  The extractor parses \texttt{.pyc} headers, recovers marshalled code objects, disassembles instruction streams, recursively visits nested code objects, and records metadata such as constants, names, variables, flags, exception tables, closures, and line-position information.  It also scans source for dynamic loading APIs including \texttt{marshal}, \texttt{importlib}, custom loaders, and direct calls to \texttt{exec} over recovered code objects.

This design addresses the input-soundness question in two ways.  Multiple dataset classes reduce dependence on one distribution channel, while artifact hashes and explicit version metadata make every extracted input reproducible.  Ambiguous cases are not discarded silently: unknown bytecode versions, unparsable artifacts, encrypted or packed payloads, and dynamic loaders whose target cannot be resolved are reported as separate categories in the evaluation.
